\section{The American}

The American\footnote{\url{https://www.imdb.com/title/tt1440728}} is a 2010 film based on the novel A Very Private Gentleman\footnote{\url{https://hagelrat.blogspot.com/2011/01/americana-very-private-gentleman-martin.html}} by Martin Booth\footnote{\url{https://en.wikipedia.org/wiki/Martin\_Booth}}. Although the film deviates from the book, it contains many eloquent and elegant passages about the lifestyle and beauty of Italy. The setting of the film is a small town in the mountainous regions of Abruzzo, a province east of Rome overlooking the Adriatic sea (and one of my ancestral homes). The town is an integral part of the landscape, following its contours in a very natural way. The book points out that international agribusiness has not yet ruined the environment, so it is particularly rich in flora and fauna, particularly the butterflies that so interest the protagonist.

A striking feature of the film is that the cast is entirely European, lacking the obligatory multicultural elements of contemporary American films. This is a way of life that the post-modernists and neo-marxists detest, seeing in it either as a brake on universal progress or as representative of the idiocy of village life. If such a place still exists, it represents a traditional culture, built up over the course of a dozen centuries. Nevertheless, on the other side the neo-pagans see in its yearly rhythm of holy days and processions to the saints only a superstitious, weak, and alien society. Its defenders are AWOL.

To emphasize the point, all the male characters in the film are aged; when that generation goes, the European traditional way of life will go. The young will either ``progress'' or ineffectively chase after ancient gods.

Jack, the American, is a maker of custom high-performance weapons used in assassinations. When an assignment goes awry in Sweden, he runs off to Italy and secludes himself in the town. His chief gives him one last job, after which Jack plans to retire. The activities of the groups are played out in secret above the regularities of ordinary life, much like the chivalrous orders of the Middle Ages. The male characters sport a similar haircut, reminiscent of Julius Caesar, thus making a connection to the ancient warrior caste.

Jack works with a sense of inner detachment, although his conversations with the village priest are meant to introduce a doubt into his mind. Nevertheless, it is interesting that Jack retains some characteristics of the Primordial State. First of all, there is his intense concentration. He is intensely focused on his task, which often involves modifications of spare parts.

The other factor is his acute sensory awareness. He is conscious of everything in his environment that could affect him and has almost a sixth sense in his anticipation of danger. This is not due to a better sense apparatus, but rather to his detachment and lack of a personal element. He has no close friends or love relationships.

With women, he enjoys before possessing. After a visit to the prostitute Carla when she reaches orgasm, he complains that he came to get pleasure not give it. But she desires to possess in order to enjoy, so she pursues a personal relationship with Jack.

His personal relationships with Carla and the priest lead him to his final act. In the Bhagavad Gita, the warrior is justified to kill his oppressor if there is no thought of reward. His career, however, was not directed against oppressors, and he was well rewarded. In the end he identifies and does kill his oppressor and, in a chivalrous act, gives up his reward to another.


\flrightit{Posted on 2011-01-18 by Cologero}

